% Build in this directory: pdflatex -interaction=nonstopmode -halt-on-error class11_handout.tex
\documentclass[a4paper,10pt,pdflatex,ja=standard,jafont=ipaex-type1,
  magstyle=real,scale=1]{bxjsarticle}
\usepackage{lmodern}
\usepackage{amsmath,amssymb,mathtools}
\usepackage{geometry}
\geometry{reset,left=22mm,right=22mm,top=20mm,bottom=17mm,
  headheight=13pt,headsep=6mm,footskip=10mm}
\usepackage{xcolor}
\usepackage{enumitem}
\usepackage{array,booktabs,tabularx}
\usepackage{fancyhdr}

\definecolor{ink}{HTML}{193A5C}
\definecolor{accent}{HTML}{237D82}
\definecolor{soft}{HTML}{EAF3F3}
\definecolor{muted}{HTML}{52616B}
\newcommand{\coursename}{解析学 II}
\newcommand{\handoutnumber}{11}
\pagestyle{fancy}
\fancyhf{}
\fancyhead[L]{\small\color{ink}\coursename{}・授業要約と演習}
\fancyhead[R]{\small\color{ink}第\handoutnumber 回}
\fancyfoot[C]{\small\color{ink}\thepage/2}
\renewcommand{\headrulewidth}{0.3pt}
\setlength{\parindent}{0pt}
\setlength{\parskip}{2pt}
\setlength{\abovedisplayskip}{5pt}
\setlength{\belowdisplayskip}{5pt}
\setlength{\abovedisplayshortskip}{3pt}
\setlength{\belowdisplayshortskip}{4pt}
\setlist[itemize]{leftmargin=1.8em,itemsep=1pt,topsep=2pt,parsep=0pt}

\newcommand{\handouttitle}[1]{%
  \begin{center}
    {\Large\color{ink}#1}\par\vspace{3mm}
    {\color{accent}\rule{0.88\linewidth}{0.7pt}}
  \end{center}\vspace{-2mm}}
\newcommand{\handout}[2]{%
  \renewcommand{\handoutnumber}{#1}%
  \handouttitle{第#1回\quad #2}}
\newcommand{\topic}[1]{\par\vspace{5pt}%
  {\large\sffamily\mdseries\color{accent}#1}\par\vspace{2pt}}
\newcommand{\keybox}[1]{\par\vspace{3pt}\begingroup
  \setlength{\fboxsep}{4pt}%
  \colorbox{soft}{\parbox{\dimexpr\linewidth-2\fboxsep\relax}{\small #1}}%
  \endgroup\par\vspace{2pt}}
\newenvironment{problems}
  {\begin{enumerate}[leftmargin=2em,itemsep=3pt,topsep=3pt,parsep=0pt,
    font=\color{accent}]}
  {\end{enumerate}}


\newcommand{\examplelabel}[1]{{\sffamily\mdseries\color{ink}#1}}
\newcommand{\answerroom}{\par\vspace{16mm}}

\begin{document}
\fontsize{10pt}{13.5pt}\selectfont
\setlength{\abovedisplayskip}{4pt}
\setlength{\belowdisplayskip}{4pt}
\handout{11}{線積分}
\topic{曲線の表示と弧長要素}
区分的に $C^1$ 級の正則曲線 $C$ を
$\vec r(t)=(x(t),y(t),z(t))$（$a\le t\le b$）で表す。
パラメータが増える方向を曲線の向きとする。
$$
 ds=|\vec r'(t)|\,dt
 =\sqrt{x'(t)^2+y'(t)^2+z'(t)^2}\,dt,\qquad
 \operatorname{Length}(C)=\int_a^b|\vec r'(t)|\,dt.
$$
積分は曲線の各滑らかな部分で計算して足し合わせる。

\topic{スカラー場の線積分：弧長に沿う総量}
曲線上の連続なスカラー場 $f$ に対して
$$
 \int_Cf\,ds:=\int_a^bf(\vec r(t))\,|\vec r'(t)|\,dt.
$$
$f$ が線密度なら、この積分は曲線状の物体の質量を表す。
弧長要素は非負であり、曲線の向きを逆にしても値は変わらない。
\par\examplelabel{例：単位円周上の $\int_Cx^2\,ds$。}
$\vec r(t)=(\cos t,\sin t)$、$0\le t\le2\pi$、$|\vec r'|=1$ なので
$$
 \int_Cx^2\,ds=\int_0^{2\pi}\cos^2t\,dt=\pi.
$$

\topic{ベクトル場の線積分：仕事と循環}
連続なベクトル場 $\vec F=(P,Q,R)$ に対して
$$
 \int_C\vec F\cdot d\vec r
 :=\int_a^b\vec F(\vec r(t))\cdot\vec r'(t)\,dt
 =\int_C(P\,dx+Q\,dy+R\,dz).
$$
曲線の向きを逆にすると符号が変わる。閉曲線上の積分は $\oint_C$ とも書く。
力の場なら仕事、速度場なら接線方向の流れの循環を表す。
\par\examplelabel{例：$\vec F=(-y,x)$ と単位円周の反時計回り。}
$$
 \vec F(\vec r(t))=(-\sin t,\cos t)=\vec r'(t),\qquad
 \oint_C\vec F\cdot d\vec r=\int_0^{2\pi}1\,dt=2\pi.
$$

\topic{ポテンシャルと線積分の基本定理}
$\phi$ が $C^1$ 級、$\vec F=\nabla\phi$、$C$ が $A$ から $B$ への曲線なら
$$
 \int_C\nabla\phi\cdot d\vec r
 =\int_a^b\frac d{dt}\phi(\vec r(t))\,dt
 =\phi(B)-\phi(A).
$$
この値は端点だけで決まり、経路によらない。特に閉曲線では0になる。
たとえば $\vec F=(x,y)$ は $\phi=(x^2+y^2)/2$ の勾配である。

\topic{再パラメータ化と向き}
同じ曲線を重複なく走る正則な表示へ変えても、弧長積分の値は変わらない。
ベクトル場の線積分は、向きを保つ表示変更で値が保存される。
曲線を何周も走れば、その回数に応じて積分も加算される。
\keybox{$f\,ds$ には速度の大きさを、$\vec F\cdot d\vec r$ には速度ベクトルを使う。始点・終点と進む向きを必ず指定する。}

\newpage
\handouttitle{第11回\quad 演習問題}
計算過程を記し、必要な条件・積分領域・向きを明示すること。
\begin{problems}
\item 半径 $R>0$ の円周 $C$ について、$\int_C1\,ds$ と $\int_Cx^2\,ds$ を求めよ。\answerroom
\item $\vec r(t)=(t,t^2)$（$0\le t\le1$）で表される曲線に沿って、$\int_Cx\,ds$ を求めよ。\answerroom
\item $\vec F=(-y,x)$ を単位円周で時計回りに積分せよ。反時計回りの値との関係も説明せよ。\answerroom
\item $\vec F=(x,y)$ の、$(0,0)$ から $(1,2)$ への線分に沿う仕事を、直接の積分とポテンシャルの両方で求めよ。\answerroom
\item $\vec F=(2xy,x^2+2y)$ のポテンシャルを求め、$(0,0)$ から $(1,2)$ への任意の区分的に滑らかな曲線に沿う線積分を求めよ。\answerroom
\item $\vec F=(-y,x)$ を $(0,0)$ から $(1,1)$ へ積分する。対角線に沿う経路と、$(1,0)$ を経由する2本の線分の経路でそれぞれ計算し、経路依存性を確認せよ。\answerroom
\end{problems}
\end{document}
